\begin{helper}[Inequality (3.4)]
Let $n\ge2$, $p>1$ and $r>0$.  Then there is a constant $c>0$, depending only on
$n$, $p$ and $r$, with the following property.

Let $\varepsilon>0$, $\gamma_0>0$ and $\gamma_1\le r\gamma_0$, let
$\Omega\subseteq\R^n$ be open, let $A$ satisfy (Q1) of the thesis--(Q2) of the thesis on
$\Omega$ with the parameters $p,\varepsilon,\gamma_0,\gamma_1$
(\noderef{EllipticStructureQ12}), assume the cited regularity result (C1) for
$\varepsilon>0$ (\noderef{CitedRegularityEpsPos}) with these parameters on
$\Omega$ for the field $A$, let $u$ be a weak solution of
$\operatorname{div}A(\nabla u)=0$ in $\Omega$ (\noderef{IsWeakSolutionOn}) with
weak gradient $g=\nabla u$ and with $u\in W^{2,2}_{\mathrm{loc}}(\Omega)$, i.e.\
each $g^k=D_ku$ lies in $W^{1,2}_{\mathrm{loc}}(\Omega)$ with weak gradient
$Dg^k$ (\noderef{MemW1pLocOn}), so that $D_jD_ku=(Dg^k)^j$.  Put
$w=\varepsilon+|\nabla u|^2$.  Then for every $\zeta\in C_0^\infty(\Omega)$,
every $k\in\{1,\dots,n\}$ and every $\ell\in\R$, in both the $+$ and the $-$
case,
\[
\int_\Omega w^{\frac p2-1}\bigl|\nabla(D_ku-\ell)^{\pm}\bigr|^2\zeta^2\,dx
\;\le\;c
\int_\Omega w^{\frac p2-1}\bigl|(D_ku-\ell)^{\pm}\bigr|^2|\nabla\zeta|^2\,dx .
\]
On the left-hand side the weak gradient of $(D_ku-\ell)^+$ is
$\nabla(D_ku)\cdot\chi_{\{D_ku>\ell\}}$, so the integral is that of
$w^{p/2-1}|\nabla D_ku|^2\zeta^2$ over $\Omega\cap\{D_ku>\ell\}$, and
symmetrically for the $-$ case over $\Omega\cap\{D_ku<\ell\}$; also
$(D_ku-\ell)^+=\max\{D_ku-\ell,0\}$, $(D_ku-\ell)^-=\max\{\ell-D_ku,0\}$ and
$|\nabla\zeta|^2=\sum_i(D_i\zeta)^2$.

\noindent\textbf{The constant, and the quantifier order.}  The thesis displays
(3.4) of the thesis with the constant $4(\gamma_1/\gamma_0)^2$ (p.~16): a constant
depending only on the structure ratio, hence on nothing that is bound after it.
The statement above records that dependence in the way the paper's Note
prescribes --- bind $n$, $p$ and a bound $r$ for the ratio, \emph{then} the
constant, \emph{then} all remaining data --- which is exactly the form the
companion inequality (3.3) of the thesis carries, both in the thesis, where its
constant is printed as $c(n,p,\gamma_1/\gamma_0)$ on the same page, and in this
development.  The structure ratio enters through the single hypothesis
$\gamma_1\le r\gamma_0$, so $c$ is a function of $n,p,r$ alone and of nothing
else; the proof exhibits $c=4n^2r^2$.  Both the $+$ and the $-$ case are
asserted with that one constant.

The ordering $\gamma_0\le\gamma_1$ is \emph{not} assumed: it is a consequence of
(Q1) of the thesis--(Q2) of the thesis, derived in the first step of the proof, and the
statement is therefore free of it.

The hypothesis (C1) is carried for the same reason as in (3.3) of the thesis: the differentiated equation
(3.2) of the thesis (\noderef{DifferentiatedEquationWeak}), which this inequality is
obtained from by substitution, carries it as an explicit hypothesis, because it
is (C1) that puts $u$ in $W^{2,2}_{\mathrm{loc}}\cap C^1$ and so makes
(3.2) of the thesis available at all.  Without it the present statement would have no
route to its own hypothesis, and the Note for the formalization permits a cited
result to be taken as an explicit hypothesis where it is used.
\end{helper}
\begin{proof}
Put
\[
c=4n^2r^2 ,
\]
which is positive because $n\ge2$ and $r>0$, and which is built from $n$ and $r$
alone, hence is chosen before
$\varepsilon,\gamma_0,\gamma_1,\Omega,A,u,g,Dg,\zeta,k,\ell$, as the quantifier
order of the statement requires.  Now fix those data satisfying all the
hypotheses.

Throughout, $k$ and $\ell$ are the index and the level of the statement, and we
write
\[
w(x)=\varepsilon+|g(x)|^2,\qquad
W(x)=w(x)^{\frac p2-1},\qquad
M_{ij}(x)=D_{h_j}A^i\bigl(g(x)\bigr),\qquad
K=\operatorname{supp}\zeta ,
\]
so that $M_{ij}(x)$ is the $i$-th component of the derivative of $A$ at $g(x)$
in the direction of the $j$-th standard basis vector, $K$ is compact and
$K\subseteq\Omega$, and $w\ge\varepsilon>0$ everywhere.  Put also
\[
X^{+}=\int_{\Omega\cap\{D_ku>\ell\}}W\,|Dg^k|^2\zeta^2\,dx,
\qquad
Y^{+}=\int_\Omega W\,\bigl((D_ku-\ell)^{+}\bigr)^2|\nabla\zeta|^2\,dx,
\]
and let $X^{-},Y^{-}$ be the same with $\{D_ku<\ell\}$ and $(\ell-D_ku)^{+}$ in
place of $\{D_ku>\ell\}$ and $(D_ku-\ell)^{+}$.  The assertion is
$X^{\pm}\le cY^{\pm}$.  All four quantities are integrals of nonnegative
functions, hence well defined in $[0,\infty]$; Step 6 shows they are finite.

The route produces the sharper bound
$X^{\pm}\le4n^2(\gamma_1/\gamma_0)^2Y^{\pm}$, and $c$ absorbs it because
$\gamma_1\le r\gamma_0$ with $\gamma_0>0$ gives $\gamma_1/\gamma_0\le r$.  The
dimensional factor $n^2$ is not slack in that bound, and the constant
$4(\gamma_1/\gamma_0)^2$ printed in the thesis is not what
(Q1) of the thesis--(Q2) of the thesis give: (Q2) of the thesis bounds the \emph{entries}
$D_{h_j}A^i$, so the bilinear form $\sum_{i,j}M_{ij}a_jb_i$ admits only the
bound $n\gamma_1W|a||b|$ of $(3)$ below, and that is attained --- a matrix all
of whose entries equal $\gamma_1W$ has quadratic form $n\gamma_1W|\xi|^2$ at
$\xi=(1,\dots,1)/\sqrt n$.  Reading (Q2) of the thesis instead as a bound on the
matrix norm of $(M_{ij})$ would give the printed constant.  Since the thesis
uses (3.4) of the thesis only through the constant of (2.1) of the thesis for $f=D_ku$,
which the lemmas of \S2 let depend on $n$ and on that constant, the existential
form above is what the statement needs and the factor costs nothing; the same
factor is visible in the thesis' own derivation of (3.3) of the thesis on the same
page, whose intermediate display carries $\gamma_1n^2$.

\emph{Step 0: the empty case, and a point of $\Omega$.}  If
$\Omega=\emptyset$ then all four of $X^{\pm}$ and $Y^{\pm}$ are integrals over
subsets of $\Omega$, hence over the empty set, hence $0$, and both assertions
read $0\le c\cdot0$; so assume $\Omega\ne\emptyset$ from now on.  Being open and
nonempty, $\Omega$ contains a ball of positive radius and therefore has positive
measure $\omega_n\rho^n>0$ (Preamble: the volume of a ball).  \emph{Every
subsequent step uses this, and two of them use it essentially}, namely Step 1 and
the measurability paragraph of Step 6, both of which read a clause of
\noderef{EllipticStructureQ12} at a single point of $\Omega$.

\emph{Step 1: the structure conditions hold at every point.}  The clauses
(Q1) of the thesis and (Q2) of the thesis of \noderef{EllipticStructureQ12} are clauses
(5) and (6) of that node and are stated for almost every $x$ --- precisely, as
``for almost every $x$, if $x\in\Omega$ then \dots'' --- so their exceptional set
is a subset of $\Omega$ and the two clauses say nothing at all when
$\Omega=\emptyset$.  That case is disposed of in Step 0, so $\Omega$ has positive
measure, and the union of the two exceptional null sets therefore cannot contain
the whole of $\Omega$: there is a single point $x\in\Omega$ at which both clauses
hold.  Fix such an $x$ and read clauses (5) and (6) there at $z=0$.  Because the
field of the present node is $(x,z,h)\mapsto A(h)$, the derivative
$D_{h_j}A^i(x,z,h)$ does not depend on $(x,z)$, so what those two clauses give at
that one point is a statement about $h$ alone and holds for every $h$.  Hence,
for every $h$ with $\varepsilon+|h|^2\ne0$ --- which by $\varepsilon>0$ is every
$h$ --- and all $\xi$, $i$, $j$,
\[
\gamma_0\bigl(\varepsilon+|h|^2\bigr)^{\frac p2-1}|\xi|^2
\le\sum_{i,j}D_{h_j}A^i(h)\,\xi^i\xi^j ,
\qquad
\bigl|D_{h_j}A^i(h)\bigr|\le\gamma_1\bigl(\varepsilon+|h|^2\bigr)^{\frac p2-1} .
\tag{1}
\]
Reading $(1)$ at $h=g(x)$: for every $x$ and every $\xi\in\R^n$,
\[
\gamma_0W(x)|\xi|^2\le\sum_{i,j}M_{ij}(x)\xi^i\xi^j ,
\qquad
|M_{ij}(x)|\le\gamma_1W(x) .
\tag{2}
\]
From the second bound and the Cauchy--Schwarz inequality applied twice,
$\sum_j|a_j|\le\sqrt n\,|a|$ and $\sum_i|b_i|\le\sqrt n\,|b|$, we get the
bilinear estimate used below:
\[
\Bigl|\sum_{i,j}M_{ij}(x)a_jb_i\Bigr|
\le\gamma_1W(x)\Bigl(\sum_j|a_j|\Bigr)\Bigl(\sum_i|b_i|\Bigr)
\le n\,\gamma_1W(x)\,|a|\,|b|
\qquad(a,b\in\R^n).
\tag{3}
\]
The same two bounds give the ordering of the structure constants, which the
statement does not assume and which Step 5 needs:
\[
0<\gamma_0\le\gamma_1 .
\tag{4}
\]
Indeed, take $h=\xi=e_1$ in $(1)$, where $e_1$ is the first standard basis
vector, so that $\varepsilon+|h|^2=\varepsilon+1$ and $|\xi|^2=1$; the first
bound gives
$\gamma_0(\varepsilon+1)^{\frac p2-1}\le\sum_{i,j}D_{h_j}A^i(e_1)e_1^ie_1^j
=D_{h_1}A^1(e_1)$, the second gives
$D_{h_1}A^1(e_1)\le|D_{h_1}A^1(e_1)|\le\gamma_1(\varepsilon+1)^{\frac p2-1}$,
and $(\varepsilon+1)^{\frac p2-1}>0$ may be cancelled.  In particular
$\gamma_1>0$, so the ratio $\gamma_1/\gamma_0$ is positive and, by the
hypothesis $\gamma_1\le r\gamma_0$, at most $r$.

\emph{Step 2: what (C1) supplies.}  Instantiate
\noderef{CitedRegularityEpsPos}, a hypothesis of the present node, at the
present $n,p,\varepsilon,\gamma_0,\gamma_1,\Omega,A$: its own hypotheses are
$0<\varepsilon$ and the bundle (Q1) of the thesis--(Q2) of the thesis
(\noderef{EllipticStructureQ12}), both hypotheses here, and it is applied to the
weak solution $u$ with weak gradient $g$ (\noderef{IsWeakSolutionOn}).  Among
its conclusions is a field $g'$ with $g'$ continuous on $\Omega$ and $g=g'$
almost everywhere on $\Omega$.  \emph{It does not assert continuity of the given
$g$, and cannot: a weak gradient is an almost-everywhere class, so continuity is
a property of a representative.  Every pointwise consequence drawn below is
therefore drawn for $g'$, and transported back to $g$ across the
almost-everywhere identity, which every integral and every almost-everywhere
assertion is blind to.}

Write $w'=\varepsilon+|g'|^2$ and $W'=(w')^{\frac p2-1}$ for the versions of $w$
and $W$ built from $g'$.  Since $g'$ is continuous on $\Omega$ and
$K=\operatorname{supp}\zeta$ is a compact subset of $\Omega$, the function
$|g'|$ is bounded on $K$, say $|g'|\le G$ there, so
\[
\varepsilon\le w'\le\varepsilon+G^2\quad\text{on }K,
\qquad\text{and therefore}\qquad
0<m_W\le W'\le M_W<\infty\quad\text{on }K,
\]
where $m_W$ and $M_W$ are the minimum and the maximum of the continuous
positive function $t\mapsto t^{\frac p2-1}$ on the compact interval
$[\varepsilon,\varepsilon+G^2]\subseteq(0,\infty)$.  Transporting back: $g=g'$
almost everywhere on $\Omega$, and $w$, $W$ are obtained from $g$ by the same
two formulas that produce $w'$, $W'$ from $g'$, so $w=w'$ and $W=W'$ almost
everywhere on $\Omega$, whence
\[
0<m_W\le W\le M_W<\infty\quad\text{almost everywhere on }K.
\tag{5}
\]
The almost-everywhere form is all that $(5)$ is ever used for: its only two uses
are the domination bounds of Step 6 and the finiteness bounds for $X^{\pm}$ and
$Y^{\pm}$ there, and a bound valid almost everywhere dominates the same
integral as one valid everywhere.  This is the second use of
(C1) in this proof: the first is that it is a hypothesis of
\noderef{DifferentiatedEquationWeak}, which Step 5 invokes.

\emph{Step 3: the test function in the $+$ case.}  We produce a pair
$(\varphi^{+},g_{\varphi^{+}})$ admissible in
\noderef{DifferentiatedEquationWeak}, namely
\[
\varphi^{+}=(D_ku-\ell)^{+}\zeta^2,
\qquad
g_{\varphi^{+}}=\zeta^2\,Z^{+}+2(D_ku-\ell)^{+}\zeta\,\nabla\zeta ,
\qquad
Z^{+}=\chi_{\{D_ku>\ell\}}\,Dg^k .
\]

\emph{(3a) An intermediate open set.}  Let $L\subseteq\Omega$ be compact.  The
distance $\operatorname{dist}(L,\R^n\setminus\Omega)$ is positive when
$\R^n\setminus\Omega\ne\emptyset$, being the distance between a compact set and
a disjoint closed set; choose $\delta>0$ with
$\delta<\operatorname{dist}(L,\R^n\setminus\Omega)$, and $\delta=1$ when
$\Omega=\R^n$.  Put $V_L=\{x:\operatorname{dist}(x,L)<\delta/2\}$.  Then $V_L$ is
open (preimage of an open half-line under a continuous function), bounded ($L$ is
bounded), contains $L$, and
$\overline{V_L}\subseteq\{x:\operatorname{dist}(x,L)\le\delta/2\}$, a closed
bounded --- hence compact --- subset of $\Omega$.  In particular
$\operatorname{volume}(V_L)<\infty$.

\emph{(3b) $D_ku$ on $V_L$.}  By hypothesis $g^k=D_ku$ lies in
$W^{1,2}_{\mathrm{loc}}(\Omega)$ with weak gradient $Dg^k$
(\noderef{MemW1pLocOn}).  Clause (4) of \noderef{WeakGradientRestrict}, applied
with $q=2$, the open sets $\Omega$ and $V_L$ and the pair $(g^k,Dg^k)$ --- its
hypotheses $1\le n$, $1\le2$, openness of both sets, $V_L\subseteq\Omega$,
compactness of $\overline{V_L}$ and $\overline{V_L}\subseteq\Omega$ all hold by
(3a) --- gives
\[
g^k\in W^{1,2}(V_L)\ \text{with weak gradient}\ Dg^k\quad(\noderef{MemW1pOn}).
\tag{6}
\]

\emph{(3c) The positive part on $V_L$.}  Apply
\noderef{WeakGradientPositivePart} with $q=2$, $c=\ell$, the open set $V_L$ and
the pair of $(6)$: its hypotheses are $1\le n$, $1\le2$, openness of $V_L$,
$\operatorname{volume}(V_L)\ne\infty$ and $(6)$.  Its conclusion is
\[
(D_ku-\ell)^{+}\in W^{1,2}(V_L)\ \text{with weak gradient}\ Z^{+} ,
\tag{7}
\]
since $(D_ku-\ell)^{+}=\max\{D_ku-\ell,0\}$ and the node's conclusion names
exactly the field that equals $Dg^k$ where $\ell<D_ku$ and $0$ elsewhere, which
is $Z^{+}$.

\emph{(3d) From the sets $V_L$ to $\Omega$.}  We claim
\[
(D_ku-\ell)^{+}\in W^{1,2}_{\mathrm{loc}}(\Omega)\ \text{with weak gradient}\ Z^{+}
\quad(\noderef{MemW1pLocOn}).
\tag{8}
\]
For the local integrability clause, let $L\subseteq\Omega$ be compact.  Then
$|(D_ku-\ell)^{+}|\le|g^k|+|\ell|$ and $|Z^{+}|\le|Dg^k|$ pointwise; by
\noderef{MemW1pLocOn} for the pair $(g^k,Dg^k)$ both $|g^k|$ and $|Dg^k|$ lie in
$L^2(L)$, the constant $|\ell|$ lies in $L^2(L)$ because $|L|<\infty$
(Preamble), and the $L^2$ triangle inequality (Minkowski, Preamble) together
with domination (Preamble) puts $(D_ku-\ell)^{+}$ and $|Z^{+}|$ in $L^2(L)$.
For the weak-gradient identity, let $\theta\in C^\infty(\R^n)$ have compact
support with $\operatorname{supp}\theta\subseteq\Omega$, and fix an index $i$.
Apply (3a)--(3c) with $L=\operatorname{supp}\theta$: then $\theta$ is a test
function for the open set $V_L$, because
$\operatorname{supp}\theta=L\subseteq V_L$, so the identity of
\noderef{HasWeakGradientOn} contained in $(7)$ gives
\[
\int_{V_L}(D_ku-\ell)^{+}D_i\theta\,dx=-\int_{V_L}(Z^{+})^i\,\theta\,dx .
\]
Both integrands vanish off $\operatorname{supp}\theta\subseteq V_L\subseteq\Omega$,
so each integral over $V_L$ equals the corresponding integral over $\Omega$, and
the identity holds with $\Omega$ in place of $V_L$.  As $\theta$ and $i$ were
arbitrary, this is $\mathrm{HasWeakGradientOn}(n,\Omega,(D_ku-\ell)^{+},Z^{+})$,
and with the local integrability just proved, $(8)$ follows.

\emph{(3e) Multiplication by $\zeta^2$.}  The function $\zeta^2$ is $C^\infty$,
being a product of two $C^\infty$ functions; it has compact support, and
$\operatorname{supp}(\zeta^2)\subseteq\operatorname{supp}\zeta=K\subseteq\Omega$.
Clause (2) of \noderef{WeakGradientSmoothMul}, applied with $q=2$, the open set
$\Omega$, the pair of $(8)$ and the factor $\zeta^2$ --- the clause that assumes
only $W^{1,2}_{\mathrm{loc}}$ and no condition on the volume of $\Omega$, which
is what is available here --- gives
\[
\varphi^{+}=(D_ku-\ell)^{+}\zeta^2\in W^{1,2}(\Omega)
\ \text{with weak gradient}\
\zeta^2Z^{+}+(D_ku-\ell)^{+}\nabla(\zeta^2)=g_{\varphi^{+}} ,
\tag{9}
\]
the last equality by the classical product rule $\nabla(\zeta^2)=2\zeta\nabla\zeta$
for the smooth function $\zeta$.  Finally $\varphi^{+}$ vanishes off $K$, so
$\operatorname{supp}\varphi^{+}\subseteq K$ is compact and contained in $\Omega$.

\emph{Step 4: the test function in the $-$ case.}  Identical, with one extra
line.  On $V_L$ the pair $(-g^k,-Dg^k)$ lies in $W^{1,2}(V_L)$ by clause (2) of
\noderef{WeakGradientLinear}, applied with $q=2$, $a=-1$, $b=0$ and both pairs
taken to be the pair of $(6)$.  Applying
\noderef{WeakGradientPositivePart} to it with $c=-\ell$ gives
\[
(\ell-D_ku)^{+}=\max\{-D_ku+\ell,0\}\in W^{1,2}(V_L)
\ \text{with weak gradient}\
Z^{-}:=-\chi_{\{D_ku<\ell\}}Dg^k ,
\]
because the node's conclusion names the field that equals $-Dg^k$ where
$-\ell<-D_ku$, i.e.\ where $D_ku<\ell$, and $0$ elsewhere.  Steps (3d) and (3e)
then apply verbatim and give
\[
\varphi^{-}=(\ell-D_ku)^{+}\zeta^2\in W^{1,2}(\Omega)
\ \text{with weak gradient}\
g_{\varphi^{-}}=\zeta^2Z^{-}+2(\ell-D_ku)^{+}\zeta\nabla\zeta ,
\tag{10}
\]
again with compact support inside $\Omega$.  Note
$(D_ku-\ell)^{-}=(\ell-D_ku)^{+}$, so $\varphi^{-}$ is the test function
$-\varphi^{\pm}$ of the thesis in the $-$ case up to the sign that the thesis
carries in front, and the identity below is unaffected by that sign because
(3.2) of the thesis is linear in the test pair.

\emph{Step 5: substitution into the differentiated equation.}  Apply
\noderef{DifferentiatedEquationWeak} with the present
$n,p,\varepsilon,\gamma_0,\gamma_1,\Omega,A,u,g,Dg$, with the index $k$ and with
the test pair $(\varphi^{+},g_{\varphi^{+}})$ of $(9)$.  Its hypotheses are
$2\le n$, $1<p$, $0<\varepsilon$, $0<\gamma_0$, the openness of $\Omega$, the
structure conditions (\noderef{EllipticStructureQ12}), the cited regularity
result (C1) (\noderef{CitedRegularityEpsPos}), the weak-solution property
(\noderef{IsWeakSolutionOn}) and the second-derivative clause
(\noderef{MemW1pLocOn}) --- all hypotheses of the present node --- together with
$\gamma_0\le\gamma_1$, which is $(4)$ of Step 1 and is the one hypothesis of
that node that the present statement does not assume, and together with
$\varphi^{+}\in W^{1,2}(\Omega)$ with weak gradient $g_{\varphi^{+}}$, compact
support and $\operatorname{supp}\varphi^{+}\subseteq\Omega$, which is Step 3.
Its conclusion is
\[
\int_\Omega\sum_{i,j}M_{ij}\,(Dg^k)_j\,(g_{\varphi^{+}})_i\,dx=0 .
\tag{11}
\]

\emph{Step 6: the two terms of $(11)$, and their finiteness.}  Write
$T_1$ and $T_2$ for the integrals obtained from $(11)$ by replacing
$g_{\varphi^{+}}$ by its two summands $\zeta^2Z^{+}$ and
$2(D_ku-\ell)^{+}\zeta\nabla\zeta$.  Both integrands vanish off $K$, because
$\zeta$ does.  On $K$ they are dominated as follows, using $(3)$ and $(5)$:
\[
\Bigl|\zeta^2\sum_{i,j}M_{ij}(Dg^k)_j(Z^{+})_i\Bigr|
\le n\gamma_1M_W\Bigl(\sup_{\R^n}\zeta^2\Bigr)|Dg^k|^2 ,
\]
\[
\Bigl|2(D_ku-\ell)^{+}\zeta\sum_{i,j}M_{ij}(Dg^k)_jD_i\zeta\Bigr|
\le n\gamma_1M_W\Bigl(\sup_{\R^n}2|\zeta||\nabla\zeta|\Bigr)
\Bigl(|Dg^k|^2+\bigl((D_ku-\ell)^{+}\bigr)^2\Bigr),
\]
where in the second line we used $|Z^{+}|\le|Dg^k|$, $(3)$ with $a=Dg^k$,
$b=\nabla\zeta$ and the elementary $2st\le s^2+t^2$.  The functions $|Dg^k|^2$
and $((D_ku-\ell)^{+})^2$ are integrable on $K$ by the local integrability
recorded in Step (3d), and $\zeta$, $\nabla\zeta$ are continuous with compact
support, hence bounded.

\emph{Measurability of the dominated functions.}  Part (a) Domination of the
Preamble's proposition on the integral requires the dominated function to be
almost everywhere strongly measurable, and a pointwise bound alone does not
supply that, so we record it.  Fix the point $x_\ast\in\Omega$ of Step 0 and read
clause (4) of \noderef{EllipticStructureQ12} --- ``for every $x\in\Omega$ and
every $z\in\R$ the map $h\mapsto A(x,z,h)$ is of class $C^1$ on
$D(\varepsilon)$'' --- at $x_\ast$ and $z=0$.  Since the present node binds
$0<\varepsilon$ \emph{strictly}, $D(\varepsilon)=\{h:\varepsilon+|h|^2\ne0\}$ is
all of $\R^n$, so $A$ is of class $C^1$ on $\R^n$ and the symbols
$D_{h_j}A^i$ used throughout are genuine partial derivatives of a $C^1$ map,
with $D_hA$ continuous on $\R^n$.  \emph{The strictness of $\varepsilon>0$ is
what makes this citation available}: at $\varepsilon\ge0$ clause (4) gives
nothing at $h=0$, and an admissible field need not be differentiable there --- for
$1<p<2$ the field $A(h)=|h|^{p-2}h$ satisfies
(Q1) of the thesis--(Q2) of the thesis at $\varepsilon=0$ while
$|A(h)-A(0)|/|h|=|h|^{p-2}\to\infty$ --- which is why the same citation is a
defect in a node that allows $\varepsilon=0$ and is sound here.  Clause (4) is
also a ``for every $x\in\Omega$'' clause, so this paragraph, like Step 1,
depends on Step 0 having produced a point of $\Omega$ at all.

By Step 2 the field $g'$ is continuous on $\Omega$ and $g=g'$ almost everywhere
on $\Omega$.  Composing the continuous $D_hA$ with $g'$, each
$M'_{ij}=D_{h_j}A^i(g'(\cdot))$ is continuous on $\Omega$, hence measurable
there; and $M_{ij}=D_{h_j}A^i(g(\cdot))$ takes the same value as $M'_{ij}$ at
every point where $g=g'$, hence almost everywhere on $\Omega$, so $M_{ij}$ is
almost everywhere equal on $\Omega$ to a continuous function and is therefore
almost everywhere strongly measurable there.  The same two sentences apply
verbatim to $W=(\varepsilon+|g|^2)^{\frac p2-1}$, whose $g'$-version
$W'=(\varepsilon+|g'|^2)^{\frac p2-1}$ is continuous on $\Omega$ --- the base
being continuous and bounded below by $\varepsilon>0$ there, and
$t\mapsto t^{p/2-1}$ being continuous on $(0,\infty)$ --- and to
$(D_ku-\ell)^{+}=\max\{g^k-\ell,0\}$, whose $g'$-version
$\max\{g'^k-\ell,0\}$ is continuous on $\Omega$.

The fields $Z^{\pm}$ are measurable rather than continuous, and this is the one
place where a \emph{set} is replaced by another, so we spell the transport out.
First, $Dg^k$ is almost everywhere strongly measurable on $\Omega$, being a weak
gradient in $W^{1,2}_{\mathrm{loc}}(\Omega)$ (\noderef{MemW1pLocOn}).  Second,
the sets $\Omega\cap\{g'^k>\ell\}$ and $\Omega\cap\{g'^k<\ell\}$ are
\emph{open}, being preimages of open half-lines under the function $g'^k$, which
is continuous on the open set $\Omega$; the corresponding assertion for $g^k$ is
false in general, since a weak gradient may be redefined on a null set, and this
is exactly the point at which the given $g$ cannot be used.  Third, let
$N=\{x\in\Omega:g(x)\ne g'(x)\}$, which is null.  Off $N$ the conditions
$g^k>\ell$ and $g'^k>\ell$ hold together or fail together, so the indicator
functions $\chi_{\Omega\cap\{g^k>\ell\}}$ and $\chi_{\Omega\cap\{g'^k>\ell\}}$
agree off $N$, that is, almost everywhere on $\Omega$.  Hence
\[
Z^{+}=\chi_{\{D_ku>\ell\}}\,Dg^k
=\chi_{\Omega\cap\{g'^k>\ell\}}\,Dg^k
\quad\text{almost everywhere on }\Omega ,
\]
and the right-hand side is the product of the indicator of an open --- hence
Borel --- set with an almost everywhere strongly measurable field, so it is
almost everywhere strongly measurable; therefore so is $Z^{+}$, being almost
everywhere equal to it.  The same argument with $<$ in place of $>$ gives
$Z^{-}$.  Finally $\zeta$ and $\nabla\zeta$ are continuous.  Every integrand appearing in
$T_1$, $T_2$, $X^{\pm}$ and $Y^{\pm}$ is therefore a finite sum of products of
these, hence almost everywhere strongly measurable on $\Omega$, which is the
hypothesis part (a) needs.

So $T_1$ and $T_2$ are finite by domination
(Preamble), $(11)$ splits as $T_1+T_2=0$, i.e.\ $T_1=-T_2$, and the same
bounds give
\[
X^{+}\le M_W\Bigl(\sup\zeta^2\Bigr)\int_K|Dg^k|^2dx<\infty,
\qquad
Y^{+}\le M_W\Bigl(\sup|\nabla\zeta|^2\Bigr)\int_K\bigl((D_ku-\ell)^{+}\bigr)^2dx<\infty .
\tag{12}
\]

\emph{Step 7: the lower bound for $T_1$.}  The integrand of $T_1$ vanishes at
every point where $Z^{+}=0$, i.e.\ off $\{D_ku>\ell\}$, and equals
$\zeta^2\sum_{i,j}M_{ij}(Dg^k)_j(Dg^k)_i$ on $\Omega\cap\{D_ku>\ell\}$, where
$Z^{+}=Dg^k$.  By the first inequality of $(2)$ with $\xi=Dg^k(x)$,
\[
T_1=\int_{\Omega\cap\{D_ku>\ell\}}\zeta^2\sum_{i,j}M_{ij}(Dg^k)_j(Dg^k)_i\,dx
\ \ge\ \gamma_0\int_{\Omega\cap\{D_ku>\ell\}}W|Dg^k|^2\zeta^2\,dx=\gamma_0X^{+},
\]
monotonicity of the integral (Preamble) being applicable because both
integrands are integrable by Step 6.

\emph{Step 8: the upper bound for $|T_2|$, and the conclusion in the $+$ case.}
The integrand of $T_2$ vanishes off $\Omega\cap\{D_ku>\ell\}$, because
$(D_ku-\ell)^{+}$ does.  By the norm inequality for the integral and $(3)$ with
$a=Dg^k$, $b=\nabla\zeta$,
\[
|T_2|\le2n\gamma_1\int_{\Omega\cap\{D_ku>\ell\}}
W\,(D_ku-\ell)^{+}\,|\zeta|\,|Dg^k|\,|\nabla\zeta|\,dx .
\]
Apply the Cauchy--Schwarz inequality (Preamble) on the set
$\Omega\cap\{D_ku>\ell\}$ with respect to the measure $W\,dx$, to the two
factors $|\zeta||Dg^k|$ and $(D_ku-\ell)^{+}|\nabla\zeta|$:
\[
|T_2|\ \le\ 2n\gamma_1
\Bigl(\int_{\Omega\cap\{D_ku>\ell\}}W|Dg^k|^2\zeta^2dx\Bigr)^{1/2}
\Bigl(\int_{\Omega\cap\{D_ku>\ell\}}W\bigl((D_ku-\ell)^{+}\bigr)^2|\nabla\zeta|^2dx\Bigr)^{1/2}
\ \le\ 2n\gamma_1\bigl(X^{+}\bigr)^{1/2}\bigl(Y^{+}\bigr)^{1/2},
\]
the last step because the first integral is $X^{+}$ and the second is at most
$Y^{+}$, the integrand being nonnegative and the domain a subset of $\Omega$.
Combining with $T_1=-T_2$ and Step 7,
\[
\gamma_0X^{+}\ \le\ T_1=|T_1|=|T_2|\ \le\ 2n\gamma_1\bigl(X^{+}\bigr)^{1/2}\bigl(Y^{+}\bigr)^{1/2} .
\tag{13}
\]
(Here $T_1\ge\gamma_0X^{+}\ge0$, so $T_1=|T_1|$.)  If $X^{+}=0$ then
$X^{+}\le cY^{+}$ holds because the right-hand side is nonnegative.  If
$X^{+}>0$ then, $X^{+}$ being finite by $(12)$, we may divide $(13)$ by
$(X^{+})^{1/2}>0$ and obtain
$\gamma_0(X^{+})^{1/2}\le2n\gamma_1(Y^{+})^{1/2}$; squaring and dividing by
$\gamma_0^2>0$,
\[
X^{+}\ \le\ 4n^2\Bigl(\frac{\gamma_1}{\gamma_0}\Bigr)^2Y^{+}
\ \le\ 4n^2r^2\,Y^{+}=cY^{+} ,
\]
the second inequality because $0<\gamma_1/\gamma_0\le r$ by Step 1 and the
hypothesis $\gamma_1\le r\gamma_0$, so that
$(\gamma_1/\gamma_0)^2\le r^2$, and because $Y^{+}\ge0$.  This is the first of
the two asserted inequalities.

\emph{Step 9: the $-$ case.}  Substituting the pair $(\varphi^{-},g_{\varphi^{-}})$
of $(10)$ into \noderef{DifferentiatedEquationWeak} --- admissible by Step 4, the
other hypotheses being the same as in Step 5 --- gives
\[
\int_\Omega\sum_{i,j}M_{ij}(Dg^k)_j\bigl(\zeta^2Z^{-}+2(\ell-D_ku)^{+}\zeta\nabla\zeta\bigr)_i dx=0 .
\]
Write $T_1^{-}$ and $T_2^{-}$ for its two terms, finite by the computation of
Step 6 with $Z^{-}$ and $(\ell-D_ku)^{+}$ in place of $Z^{+}$ and
$(D_ku-\ell)^{+}$, so that $T_1^{-}=-T_2^{-}$.  The integrand of $T_1^{-}$
vanishes off $\{D_ku<\ell\}$ and equals
$-\zeta^2\sum_{i,j}M_{ij}(Dg^k)_j(Dg^k)_i$ there, since $Z^{-}=-Dg^k$ on that
set; hence, by the first inequality of $(2)$,
\[
-T_1^{-}=\int_{\Omega\cap\{D_ku<\ell\}}\zeta^2\sum_{i,j}M_{ij}(Dg^k)_j(Dg^k)_i\,dx
\ \ge\ \gamma_0X^{-} .
\]
The integrand of $T_2^{-}$ vanishes off $\Omega\cap\{D_ku<\ell\}$, and exactly
as in Step 8,
\[
|T_2^{-}|\le2n\gamma_1\bigl(X^{-}\bigr)^{1/2}\bigl(Y^{-}\bigr)^{1/2} .
\]
Since $-T_1^{-}=T_2^{-}$, we get
$\gamma_0X^{-}\le|T_2^{-}|\le2n\gamma_1(X^{-})^{1/2}(Y^{-})^{1/2}$, and the
dichotomy of Step 8, followed by the same passage from $\gamma_1/\gamma_0$ to
$r$, gives
$X^{-}\le4n^2(\gamma_1/\gamma_0)^2Y^{-}\le cY^{-}$.  This is the second
asserted inequality, with the same constant $c$ as the first, and the proof is
complete.
\end{proof}
